% TOP_PROBLEM: {"id": 30007441, "title": "Polynomial switch-chain mixing for all degree sequences", "category_id": 15, "category": {"id": 15, "name": "computer_science", "display_name": "Computer Science", "description": "Computational complexity, algorithms, and theoretical CS.", "slug": "computer-science", "order_index": 15, "created_at": "2026-07-31T15:26:25.671Z"}, "status": "open"}
% FIRST_POSTED: 2026-10-06
% ATTEMPT_STATUS: solved
% OPENAI_CLAIM: {"schema_version":1,"source_repo":"https://github.com/openai/math","source_commit":"adc7f1241b42e322a6451854ab7e4b4c146bf78a","release_date":"2026-10-06","match":"full","resolution":"proved","independently_reviewed":false,"second_pass":"agreed","families":[{"family":"131","title":"Rapid mixing of graph switches for every degree sequence","subject":"Theoretical computer science","manuscripts":[{"title":"Polynomial mixing of the switch chain for every graphical degree sequence","pdf_path":"preprints/Polynomial-Mixing-of-the-Switch-Chain-for-Every-Graphical-Degree-Sequence-September-25-2026/main.pdf","readme_path":"preprints/Polynomial-Mixing-of-the-Switch-Chain-for-Every-Graphical-Degree-Sequence-September-25-2026/README.md","date":"2026-09-25","theorem_ref":"preprints/Polynomial-Mixing-of-the-Switch-Chain-for-Every-Graphical-Degree-Sequence-September-25-2026/build/sections/introduction.tex#thm:main"}],"lean":{"doc_path":"lean/docs/131.md","comparators":["lean/ComparatorChallenges/SwitchChain.json","lean/ComparatorChallenges/SwitchChain.lean"],"declarations":[],"covers_main_theorem":true},"reasoning_trace":null}]}
\documentclass[11pt]{article}
\usepackage[margin=25mm]{geometry}
\usepackage{fontspec}
\setmainfont{Latin Modern Roman}
\usepackage{amsmath,amssymb,mathtools}
\usepackage{unicode-math}
\setmathfont{Latin Modern Math}
\usepackage{enumitem}
\usepackage{xurl,hyperref,bookmark}
\hypersetup{colorlinks=true,urlcolor=blue}
\setcounter{secnumdepth}{0}
\setlength{\parindent}{0pt}
\setlength{\parskip}{5pt}
\setlength{\emergencystretch}{3em}
\begin{document}
\section*{Polynomial switch-chain mixing for all degree sequences}

\subsection{OpenAI's claim}
OpenAI presents the following claims relevant to this question, with manuscript numbers following paper order in Sources. In manuscript 1 (family 131; 2026-09-25), the asserted result is that the lazy simple-undirected switch chain mixes for every graphical degree sequence in at most $2n^8$ steps at distance $1/4$, with an exact expected-polynomial sampler.

\subsection{Outline of the argument}
For manuscript 1 (family 131; 2026-09-25), pair resampling operators are compared through triple and disjoint-pair variance estimates. A global operator inequality gives the spectral gap, and decomposition into individual switches transfers it to the specified lazy chain; exact residual correction yields uniform sampling.

\subsection{Scope relative to this problem}
The result supplies the conclusion in this problem with its displayed hypotheses and quantifiers. The broader statements and applications in the manuscript are kept separate from the particular question above. The supporting result is manuscript 1 (family 131; 2026-09-25).

\subsection{Formal verification}
For manuscript 1 (family 131; 2026-09-25), the documented Lean scope is connectivity, spectral-gap control, and the stated lazy-switch total-variation bound for graphical degree sequences; the exact uniform sampler is not part of those selected estimates.

The selected statement includes the clause used for this problem.

The linked family Comparator checks are \texttt{SwitchChain.lean}.

Their named declarations include \texttt{OAI.Problem315.switch\_chain\_main}, \texttt{OAI.Problem315.switch\_chain\_tv\_bound}, \texttt{OAI.Problem315.switch\_connectivity}.

This site has not executed the Lean code. Described statement coverage is not an independent proof check.

\subsection{Status caveat}
The manuscripts were produced by an internal OpenAI model. OpenAI cautions that results without formalization can contain mistakes. This site has not independently checked the proofs; the label records the scope of OpenAI's claim.

\subsection{Sources}
\begin{itemize}
\item OpenAI, \emph{Polynomial mixing of the switch chain for every graphical degree sequence}, 2026-09-25; citation key \texttt{OAI:Polynomial-Mixing-of-the-Switch-Chain-for-Every-Graphical-Degree-Sequence-September-25-2026}. \href{https://github.com/openai/math/blob/adc7f1241b42e322a6451854ab7e4b4c146bf78a/preprints/Polynomial-Mixing-of-the-Switch-Chain-for-Every-Graphical-Degree-Sequence-September-25-2026/main.pdf}{PDF (GitHub)}; \href{https://github.com/openai/math/raw/adc7f1241b42e322a6451854ab7e4b4c146bf78a/preprints/Polynomial-Mixing-of-the-Switch-Chain-for-Every-Graphical-Degree-Sequence-September-25-2026/main.pdf}{PDF (direct)}; \href{https://github.com/openai/math/blob/adc7f1241b42e322a6451854ab7e4b4c146bf78a/preprints/Polynomial-Mixing-of-the-Switch-Chain-for-Every-Graphical-Degree-Sequence-September-25-2026/README.md}{Manuscript page and citation}; \href{https://github.com/openai/math/blob/main/preprints/Polynomial-Mixing-of-the-Switch-Chain-for-Every-Graphical-Degree-Sequence-September-25-2026/main.pdf}{Latest version}.
\item Family 131: \href{https://github.com/openai/math/blob/adc7f1241b42e322a6451854ab7e4b4c146bf78a/lean/docs/131.md}{Lean scope}.
\item \href{https://github.com/openai/math/blob/adc7f1241b42e322a6451854ab7e4b4c146bf78a/lean/ComparatorChallenges/SwitchChain.json}{Comparator SwitchChain.json}.
\item \href{https://github.com/openai/math/blob/adc7f1241b42e322a6451854ab7e4b4c146bf78a/lean/ComparatorChallenges/SwitchChain.lean}{Comparator SwitchChain.lean}.
\item \href{https://github.com/openai/math/blob/adc7f1241b42e322a6451854ab7e4b4c146bf78a/CONTENTS.md}{OpenAI catalogue entry}.
\end{itemize}
\end{document}
